% 383_Ein_Anker_ohne_v_En_ch.tex
% Johann Pascher, 29 Sept. 2026
% Masses and Planck scale without v: one chain, one anchor

\providecommand{\BM}{\textbf{[B]}}
\providecommand{\KM}{\textbf{[K]}}
\providecommand{\SM}{\textbf{[S]}}

\chapter*{Masses and Planck Scale without \texorpdfstring{$v$}{v}}
\addcontentsline{toc}{chapter}{Masses and Planck scale without v}

\begin{abstract}
The lepton ladder of the corpus writes the masses as $m_i=r_i\,\xi^{p_i}\,v$ and uses
the electroweak scale $v$ as an intermediate quantity; the Planck scale is reached
separately through the formula for $G$ with the conversion factor $C_{\text{conv}}$.
Read this way, the calculation seems to need two anchors, one for $v$ and one for
$E_P$. It does not. Doc.~149 links $v$ and $E_P$ directly:
$v=E_P/(f^4\cdot(\pi/2)\cdot10)$, i.e.\ $v/E_P=\xi^4/(5\pi)=2.012\times10^{-17}$,
$-0.23\,\%$ against the measured value \KM. Inserted into the ladder, $v$ drops out:
$m_i=\frac{r_i}{5\pi}\,\xi^{p_i+4}\,E_P$ \BM. All three lepton masses, $E_P$, $G$,
$\ell_P$ and $L_0$ then hang on a single chain, and one measured value as anchor fixes
all numbers. With $m_e$ as anchor one obtains $E_P=1.237\times10^{19}$\,GeV
($+1.34\,\%$), $G=6.50\times10^{-11}$ ($-2.6\,\%$), $\ell_P=1.595\times10^{-35}$\,m
and $L_0=2.127\times10^{-39}$\,m \KM. The deviations are the known bare residues of
the ladder, not a second anchor.
\end{abstract}

% ============================================================
\section*{Status markers}
% ============================================================
\begin{center}
\begin{tabular}{@{}lp{10cm}@{}}
\textbf{[B]} & Proven algebraically (rearrangement without additional assumption). \\[3pt]
\textbf{[K]} & Confirmed numerically (check script). \\[3pt]
\textbf{[S]} & Speculative: a hint or conjecture, not yet derived or checked in the corpus. \\
\end{tabular}
\end{center}
Measured values (CODATA 2022, PDG) serve as anchor or as comparison values. $\hbar$,
$c$ and $e$ are exact by definition in the SI; they fix only the units and are not an
additional anchor.

% ============================================================
\section{Starting point}
% ============================================================
Two routes stand side by side in the corpus:
\begin{enumerate}
\item \textbf{Lepton ladder} (Docs.~006, 352): $m_i=r_i\,\xi^{p_i}\,v$ with
  $r=(4/3,\;16/5,\;25/9)$ and $p=(3/2,\;1,\;2/3)$ for $e,\mu,\tau$. The mass ratios
  $m_\mu/m_e=\tfrac{12}{5}\xi^{-1/2}=207.85$ and
  $m_\tau/m_\mu=\tfrac{125}{144}\xi^{-1/3}=16.99$ do not contain $v$.
\item \textbf{Gravitation} (Docs.~012, 013, 180):
  $G=\xi^2/(4m_e)\cdot C_{\text{conv}}\cdot K_{\text{frak}}$, and from it $E_P$ and $\ell_P$.
\end{enumerate}
Taking $v$ from $m_e$ and $E_P$ from the $G$ formula uses two measured values: $m_e$
and, through $C_{\text{conv}}$, the measured value of $G$. This is how Doc.~375
calculates. A chain meant to derive all quantities from $\xi$ alone, however, may have
only one anchor; all other numbers must follow. Doc.~149 supplies the link.

% ============================================================
\section{The relation between \texorpdfstring{$v$}{v} and \texorpdfstring{$E_P$}{E\_P}}
% ============================================================
Doc.~149 derives the electroweak scale from the Planck energy: dilution over the $f^4$
cells of the four-dimensional lattice, projection onto the three-dimensional half-space
($\pi/2$), and a factor $1/10$ for the step to the electroweak scale:
\begin{equation}
  v=\frac{E_P}{f^4\cdot(\pi/2)\cdot10}
  \qquad\Longleftrightarrow\qquad
  \frac{v}{E_P}=\frac{\xi^4}{5\pi}\quad\BM
\end{equation}
with $f=1/\xi$. The rearrangement is exact. Doc.~149 labels the factor $10$ a motivated
choice; it is not derived \SM.

\textbf{Numerical value.} $v/E_P=\xi^4/(5\pi)=2.0120\times10^{-17}$. The measured value
is $v/E_P=246.22\,\text{GeV}/1.22089\times10^{19}\,\text{GeV}=2.0167\times10^{-17}$;
the deviation is $-0.23\,\%$ \KM. Since $v/E_P$ is dimensionless, this number needs no
anchor.

\textbf{On the value in Doc.~149.} There the calculation uses $f=7491.91$ and gives
$v=246.71$\,GeV. With $f=1/\xi=7500$ the same formula gives $v=245.65$\,GeV \KM. Both
values are calculated correctly; they differ in the $f$ inserted.

% ============================================================
\section{The masses without \texorpdfstring{$v$}{v}}
% ============================================================
Inserting $v=\xi^4E_P/(5\pi)$ into the ladder gives
\begin{equation}
  m_i=\frac{r_i}{5\pi}\;\xi^{\,p_i+4}\;E_P\quad\BM
\end{equation}
With the measured value $E_P=1.22089\times10^{22}$\,MeV for comparison:

\begin{center}
\begin{tabular}{@{}lllrr@{}}
\toprule
\textbf{Lepton} & \textbf{Formula} & \textbf{Coefficient} & \textbf{Value} & \textbf{vs.\ measurement}\\
\midrule
$e$    & $\frac{4}{15\pi}\,\xi^{11/2}E_P$ & $0.08488$ & $0.5043$\,MeV  & $-1.32\,\%$\\[3pt]
$\mu$  & $\frac{16}{25\pi}\,\xi^{5}E_P$   & $0.20372$ & $104.81$\,MeV  & $-0.80\,\%$\\[3pt]
$\tau$ & $\frac{5}{9\pi}\,\xi^{14/3}E_P$ & $0.17684$ & $1780.9$\,MeV  & $+0.23\,\%$\\
\bottomrule
\end{tabular}
\end{center}
\KM\ The mass ratios are the same as in the ladder, since $E_P$ and $5\pi$ cancel. The
deviations combine the bare residue of the ladder against $v$ (Doc.~352) and the
$-0.23\,\%$ of the relation $v/E_P$.

% ============================================================
\section{One chain, one anchor}
% ============================================================
The chain reads:
\begin{equation}
  \xi\;\longrightarrow\;\frac{m_i}{E_P}=\frac{r_i}{5\pi}\,\xi^{p_i+4}
  \;\longrightarrow\;E_P\;\longrightarrow\;
  G=\frac{\hbar c^5}{E_P^2},\quad
  \ell_P=\frac{\hbar c}{E_P},\quad
  L_0=\xi\,\ell_P .
\end{equation}
Apart from the units, everything is fixed by $\xi$. A single measured value fixes the
numbers. With $m_e=0.51099895$\,MeV as anchor:

\begin{center}
\begin{tabular}{@{}lrr@{}}
\toprule
\textbf{Quantity} & \textbf{from anchor $m_e$} & \textbf{vs.\ measurement}\\
\midrule
$E_P$    & $1.2372\times10^{19}$\,GeV & $+1.34\,\%$\\
$G$      & $6.4995\times10^{-11}$\,m$^3$\,kg$^{-1}$\,s$^{-2}$ & $-2.6\,\%$\\
$\ell_P$ & $1.5950\times10^{-35}$\,m & $-1.32\,\%$\\
$L_0=\xi\ell_P$ & $2.1266\times10^{-39}$\,m & $-1.32\,\%$\\
\bottomrule
\end{tabular}
\end{center}
\KM\ Since $G\propto E_P^{-2}$, the relative deviation of $G$ is twice that of $E_P$.

\textbf{Choice of anchor.} The chain can start at any point; which quantity serves as
anchor is a question of accuracy, not of derivation. The deviations of the other
quantities then depend on the bare residue of the chosen entry point:

\begin{center}
\begin{tabular}{@{}lrrr@{}}
\toprule
\textbf{Anchor} & $E_P$ & $G$ & $\ell_P$, $L_0$\\
\midrule
$m_e$    & $+1.34\,\%$ & $-2.62\,\%$ & $-1.32\,\%$\\
$m_\mu$  & $+0.81\,\%$ & $-1.60\,\%$ & $-0.80\,\%$\\
$m_\tau$ & $-0.23\,\%$ & $+0.45\,\%$ & $+0.23\,\%$\\
$v$      & $+0.23\,\%$ & $-0.46\,\%$ & $-0.23\,\%$\\
\bottomrule
\end{tabular}
\end{center}
\KM\ In every case exactly one measured value is used. The spread between the rows is
the bare residue of the ladder; it shrinks to the extent that Doc.~352 resolves this
residue.

% ============================================================
\section{Relation to the \texorpdfstring{$G$}{G} formula with \texorpdfstring{$C_{\text{conv}}$}{C\_conv}}
% ============================================================
The formula $G=\xi^2/(4m_e)\cdot C_{\text{conv}}\cdot K_{\text{frak}}$ from Doc.~012
gives $6.6745\times10^{-11}$, i.e.\ $+0.003\,\%$. This does not contradict the
$-2.6\,\%$ above: $C_{\text{conv}}=7.783\times10^{-3}$ is chosen so that the SI
conversion matches the measured value of $G$. In this form $m_e$ and $G$ both enter as
anchors. The one-anchor chain with $m_e$ corresponds to
$C_{\text{conv}}=7.58\times10^{-3}$ \KM; the difference of $2.6\,\%$ is exactly the
residue of the chain. $C_{\text{conv}}$ remains the SI conversion of $G$; the chain only
shows that it is not needed as a second measured value if the residue is accepted.

Doc.~375 calculates $v/E_P$ along this route ($v$ from $m_e$, $E_P$ via
$C_{\text{conv}}$) and obtains $2.0389\times10^{-17}$. The direct value
$\xi^4/(5\pi)=2.0120\times10^{-17}$ lies lower by the factor $1.0134$; this is exactly
the residue $-1.32\,\%$ of the electron mass in the chain \KM. Both routes are therefore
consistent.

% ============================================================
\section{Notes in other documents}
% ============================================================
\begin{itemize}
\item \textbf{Doc.~149:} note at the formula for $v$ that it gives $245.65$\,GeV with
  $f=7500$ and means $v/E_P=\xi^4/(5\pi)$.
\item \textbf{Doc.~180:} note at $\ell_P$ that the value $1.616\times10^{-35}$\,m is the
  measured value used as anchor and that the one-anchor chain gives
  $1.595\times10^{-35}$\,m.
\item \textbf{Doc.~375:} note that $v/E_P$ follows directly as $\xi^4/(5\pi)$, without
  $C_{\text{conv}}$ and without an anchor.
\end{itemize}

% ============================================================
\section{Result}
% ============================================================
In the mass calculation $v$ is an intermediate quantity and can be eliminated:
$m_i=\frac{r_i}{5\pi}\xi^{p_i+4}E_P$ \BM. The lepton masses, $E_P$, $G$, $\ell_P$ and
$L_0$ thus hang on one chain that is fixed by $\xi$ apart from the units. A single
measured value suffices as anchor; which one is free. The deviations of $0.2$ to
$1.3\,\%$ in $E_P$ are the bare residues of the ladder and of the relation $v/E_P$
($-0.23\,\%$), not the consequence of a missing anchor. The factor $10$ in this
relation is the open point of the chain \SM.

\medskip\noindent Check script: \texttt{2/python/Dok383\_Skripte/pruef\_383\_ein\_anker.py}, 20/20.

\begin{thebibliography}{9}
\bibitem{CODATA2022} E.~Tiesinga et al., \emph{CODATA Recommended Values of the
  Fundamental Physical Constants: 2022}, NIST (2024).
\bibitem{PDG2024} Particle Data Group (S.~Navas et al.), \emph{Review of Particle
  Physics}, Phys.\ Rev.\ D 110 (2024) 030001.
\end{thebibliography}
